\chapter{Ideomotor Theory and the Predictive Brain}
\label{ch:ideomotor}

\epigraph{%
  Every representation of a movement awakens in some degree
  the actual movement which is its object.
}{--- William James, 1890}

\section{Historical Ideomotor Theory}

The ideomotor principle, articulated by William James in 1890 \cite{James1890}, states
that the mere \emph{idea} of a movement tends to produce that movement. Modern
reformulations by Greenwald (1970) and Hommel et al.\ (2001) have transformed this
philosophical observation into a testable scientific framework.

\begin{definition}{The Theory of Event Coding (TEC)}{tec-def}
  The Theory of Event Coding \cite{Hommel2001} proposes that:
  \begin{enumerate}
    \item Perception and action share a \textbf{common representational format}
          (event codes or feature codes)
    \item Actions are planned in terms of their \textbf{desired sensory effects},
          not in terms of muscle commands
    \item The mapping from desired effects to motor commands is learned through
          experience (bidirectional action-effect associations)
  \end{enumerate}
\end{definition}

\section{The Predictive Processing Framework}

Modern neuroscience interprets the brain as a \textbf{prediction machine}
\cite{Clark2013, Hohwy2013}:

\begin{insight}{The Brain as a Prediction Machine}{prediction-machine}
  The brain continuously generates predictions about incoming sensory data.
  Neural processing is dominated by \textbf{prediction error minimization}:
  \begin{equation}
    \varepsilon_k = \bm{y}_k - \hat{\bm{y}}_k = \bm{y}_k - h(\hat{\state}_k),
    \label{eq:ch5:prediction-error}
  \end{equation}
  where $\hat{\bm{y}}_k$ is the predicted observation and $\bm{y}_k$ is the actual
  observation. This error drives updates to the internal model.

  In the language of Volume~I, Chapter~6: this is exactly the
  \textbf{Kalman filter innovation}. The brain implements Bayesian state estimation.
\end{insight}

\subsection{Active Inference}

Karl Friston's active inference framework \cite{Friston2010} extends predictive
processing to action: instead of only updating beliefs to match observations,
the agent can also \emph{act} to make observations match predictions:
\begin{equation}
  \control^* = \argmin_{\control} \mathbb{E}\left[
  \sum_{k=0}^{T} \norm{\bm{y}_k - \hat{\bm{y}}_k}^2_{R^{-1}}
  + \norm{\control_k}^2_Q
  \right],
  \label{eq:ch5:active-inference}
\end{equation}
which is formally equivalent to model predictive control (Volume~I, Chapter~5).

\section{Connection to Control Theory}

\begin{center}
\renewcommand{\arraystretch}{1.3}
\begin{tabular}{@{}ll@{}}
\toprule
\textbf{Neuroscience Term} & \textbf{Control Theory Term} \\
\midrule
Forward model & State transition matrix / predictor \\
Prediction error & Innovation (Kalman filter) \\
Efference copy & Feedforward control signal \\
Active inference & Model predictive control \\
Prior beliefs & State constraints / initial conditions \\
Precision weighting & Kalman gain / observation weights \\
Free energy & Variational cost functional \\
\bottomrule
\end{tabular}
\end{center}

\section*{Exercises}

\begin{enumerate}
  \item \textbf{Conceptual.} Explain how the ideomotor principle relates to
        trajectory-centric control (Volume~II). When a golfer imagines impact,
        what ``prediction'' are they generating?

  \item \textbf{Kalman Connection.} Write out the Kalman filter equations and identify
        which term corresponds to the prediction error in the predictive processing
        framework.

  \item \textbf{(Programming.)} Implement an active inference agent that controls a
        pendulum by minimizing prediction error rather than tracking a reference.
\end{enumerate}
